\documentclass[10pt,a4paper]{amsart}
\usepackage[margin=19mm]{geometry}
\usepackage{amsmath,amssymb,mathtools}
\usepackage[scr=rsfs,cal=euler]{mathalfa}
\usepackage{xfrac}
\usepackage[unicode,hidelinks,bookmarks=false]{hyperref}
\newcommand{\ie}{i.e.\ }
\newcommand{\cf}{cf.\ }
\newcommand{\resp}{respectively\ }
\newcommand{\Subsection}{\subsection}
\providecommand{\nobreakdash}{\nobreak\hbox{-}}
\setlength{\emergencystretch}{2em}
\multlinegap0pt
\usepackage{bm}
\usepackage{bbm}
\usepackage{dsfont}
\usepackage{xspace}
\usepackage{cancel}
\usepackage{mathtools}
\usepackage{amsthm}
\makeatletter
\@ifundefined{theorem}{\newtheorem{theorem}{Theorem}}{}
\@ifundefined{corollary}{\newtheorem{corollary}[theorem]{Corollary}}{}
\@ifundefined{lemma}{\newtheorem{lemma}[theorem]{Lemma}}{}
\makeatother
\title[A smooth Alpert testing characterization of convolution type]{A smooth Alpert testing characterization of convolution type for the Fourier extension conjecture on paraboloids}
\author{Cristian Rios, Eric T. Sawyer}
\date{Source: arXiv:2512.24990v9 (CC BY 4.0)}
\begin{document}
\maketitle

\noindent\emph{Source and licence.} A smooth Alpert testing characterization of convolution type for the Fourier extension conjecture on paraboloids. Version arXiv:2512.24990v9 (CC BY 4.0), distributed under the Creative Commons Attribution 4.0 International licence, as recorded in the arXiv licence metadata for this exact version and shown on the abstract page https://arxiv.org/abs/2512.24990v9. Full rights notice: licenses/arxiv-2512.24990-CC-BY-4.0.txt.\par\smallskip

\noindent\textit{Editorial scope.} Counted unit: the Corollary whose source label is \texttt{Cor LSST} in the article (source0.tex lines 2202--2213, source label \texttt{Cor LSST}), delivered together with the article's own complete original proof of it (source0.tex lines 2215--2427, one \texttt{proof} environment of 213 lines). No mathematics is written, repaired or re-derived: statement and proof are the article's own text line for line. Required context retained uncounted: special (lines 595--598); lem special (lines 636--645); lin d (lines 1004--1010). Every definition, display and label of those lines is retained unchanged. Omitted from this package and not redistributed: 1--594, 599--635, 646--1003, 1011--2201, 2214--2214, 2428--2820. Nothing inside a retained excerpt is omitted or altered: every retained body is delivered exactly as the article's lines are written, so no omission receipt of any kind accompanies this package. Citations: the retained excerpts carry no citation command at all, so no bibliography entry of the article is transcribed and no reference is redistributed. Reference adaptation: none was needed and none was made; every reference inside the delivered unit resolves inside it, so no unresolved reference or citation is produced. Printed slips kept literal: no printed word, symbol, hypothesis or number of the counted statement or of its proof was altered. Layout and class: the article's own class is \texttt{amsart} with packages including amssymb, amsfonts, geometry, graphicx, amsmath; the wrapper is the lane's pinned \texttt{amsart} editorial wrapper at 10pt on A4, which restates the theorem-like environments the retained text needs and adds the article's own macro definitions that the retained text needs (none), with extra wrapper packages (bm, bbm, dsfont, xspace, cancel, mathtools). The delivered numbering is the unit's own. Assets: no retained span contains a figure environment, a graphics-inclusion command, a PDF-page inclusion, a table or a TikZ picture; no image file is copied into this package, the package declares no asset dependency and no image file is redistributed with it. Licence: version arXiv:2512.24990v9 of this article is distributed under the Creative Commons Attribution 4.0 International licence, as recorded in the arXiv licence metadata for this exact version and shown on the abstract page https://arxiv.org/abs/2512.24990v9. A full rights notice is delivered at licenses/arxiv-2512.24990-CC-BY-4.0.txt. Characters: every retained excerpt is pure ASCII, so the delivered engine, XeLaTeX with its default Latin Modern text font, reports no missing character.
\begin{equation}
\ell\left(  I\right)  \leq1\ \text{and }\operatorname*{dist}\left(
J,I\right)  >2^{n}\text{ and }\ell\left(  J\right)  \leq2^{n},\label{special}%
\end{equation}

\begin{lemma}
\label{lem special}For $n\in\mathbb{N}$, and $I$ and $J$ as in (\ref{special}%
), we have%
\[
\left\vert \left\langle Th_{J;\kappa},h_{I;\kappa}\right\rangle \right\vert
\leq C_{d-1}\frac{1}{\eta^{\kappa}}\left(  \frac{\ell\left(  I\right)  }%
{2^{n}}\right)  ^{\kappa+\frac{d-1}{2}}.
\]

\end{lemma}

\begin{equation}
\left\Vert \mathbb{E}_{\mathbb{G}}\mathcal{E}M_{\psi}^{\mathcal{G}}%
\mathsf{Q}_{s}^{\eta,\mathcal{G}}f\right\Vert _{L^{q}\left(  B\left(
0,2^{r}\right)  \right)  }\leq C_{\varepsilon^{\prime},\delta,\kappa
,q}2^{s\varepsilon^{\prime}}\left\Vert f\right\Vert _{L^{\infty}\left(
U\right)  },\label{lin d}%
\end{equation}

\begin{corollary}
\label{Cor LSST}Let $\frac{2d}{d-1}\leq q\leq p<\infty$. Then the inequality
(\ref{lin d}) holds if and only if%
\begin{equation}
\left\Vert \mathbb{E}_{\mathbb{G}}\mathcal{E}\widetilde{\mathsf{Q}}_{s}%
^{\eta,\mathcal{G}}f\right\Vert _{L^{q}\left(  B\left(  0,2^{\frac{s}%
{1-\delta}}\right)  \right)  }\leq C_{\varepsilon^{\prime},\delta,\kappa
,p,q}2^{s\varepsilon^{\prime}}\left\Vert f\right\Vert _{L^{p}\left(  U\right)
},\ \ \ \ \ \text{for all }f\in L^{p}\left(  U\right)  ,\label{lin d'}%
\end{equation}
and so is equivalent to the Fourier extension conjecture.
\end{corollary}

\begin{proof}
Since
\begin{align*}
\widetilde{\mathsf{Q}}_{s}^{\eta,\mathcal{G}}\mathbf{1}_{U}f  & =\mathsf{Q}%
_{s}^{\eta,\mathcal{G}}T_{\mathcal{G}}^{\eta}\mathbf{1}_{U}f=\mathsf{Q}%
_{s}^{\eta,\mathcal{G}}\left(  \mathbf{1}_{U}+\sum_{n=1}^{\infty}%
\mathbf{1}_{2^{n}U\setminus2^{n-1}U}\right)  T_{\mathcal{G}}^{\eta}%
\mathbf{1}_{U}f=\sum_{n=0}^{\infty}\mathsf{Q}_{s}^{\eta,\mathcal{G}}f_{n}\\
\text{where }f_{0}  & \equiv\mathbf{1}_{U}T_{\mathcal{G}}^{\eta}\mathbf{1}%
_{U}f\text{ and }f_{n}\equiv\mathbf{1}_{2^{n}U\setminus2^{n-1}U}%
T_{\mathcal{G}}^{\eta}\mathbf{1}_{U}f\text{ for }n\geq1,
\end{align*}
we have by (\ref{lin d}) that%
\begin{align*}
\left\Vert \mathbb{E}_{\mathbb{G}}\mathcal{E}\widetilde{\mathsf{Q}}_{s}%
^{\eta,\mathcal{G}}f\right\Vert _{L^{q}\left(  B\left(  0,2^{\frac{s}%
{1-\delta}}\right)  \right)  }  & \leq\sum_{n=0}^{\infty}\left\Vert
\mathbb{E}_{\mathbb{G}}\mathcal{E}\mathsf{Q}_{s}^{\eta,\mathcal{G}}%
f_{n}\right\Vert _{L^{q}\left(  B\left(  0,2^{\frac{s}{1-\delta}}\right)
\right)  }\\
& \leq C_{\varepsilon^{\prime},\delta,\kappa,p,q}2^{s\varepsilon^{\prime}%
}\left\Vert f_{0}\right\Vert _{L^{p}\left(  U\right)  }+\sum_{n=0}^{\infty
}A^{n}\left\Vert f_{n}\right\Vert _{L^{p}\left(  2^{n}U\right)  }%
\end{align*}
where the constant $A^{n}$ arises from the rescaling of $U$ to $2^{n}U$. From
the invertibility of $T_{\mathcal{G}}^{\eta}$ on $L^{p}$, and the rapid decay
of $f_{n}=\mathbf{1}_{2^{n}U\setminus2^{n-1}U}T_{\mathcal{G}}^{\eta}%
\mathbf{1}_{U}f$ given in the well localized properties of $T_{\mathcal{G}%
}^{\eta}$ in Lemma \ref{lem special}, we then claim that%
\[
\left\Vert \mathbb{E}_{\mathbb{G}}\mathcal{E}\widetilde{\mathsf{Q}}_{s}%
^{\eta,\mathcal{G}}f\right\Vert _{L^{q}\left(  B\left(  0,2^{\frac{s}%
{1-\delta}}\right)  \right)  }\leq C_{\varepsilon^{\prime},\delta,\kappa
,p,q}2^{s\varepsilon^{\prime}}\left\Vert f\right\Vert _{L^{p}\left(  U\right)
}\ .
\]


Indeed, it suffices to prove that there is $\tau>A$ such that
\[
\left\Vert f_{n}\right\Vert _{L^{p}\left(  U\right)  }\leq C_{\tau}2^{-\tau
n}\left\Vert f\right\Vert _{L^{p}\left(  U\right)  },\ \ \ \ \ \text{for
}n=0,1,2,...
\]
which by duality is equivalent to
\[
\left\vert \left\langle T_{\mathcal{G}}^{\eta}\mathbf{1}_{U}f,\mathbf{1}%
_{2^{n}U\setminus2^{n-1}U}g\right\rangle \right\vert =\left\vert \left\langle
f_{n},g\right\rangle \right\vert \leq C_{\tau}2^{-\tau n}\left\Vert
f\right\Vert _{L^{p}\left(  U\right)  }\left\Vert g\right\Vert _{L^{p^{\prime
}}\left(  U\right)  }.
\]
Now if $f$ is supported in $U$ and has vanishing moments up to order $\kappa$,
then%
\[
f=\sum_{I\in\mathcal{G}\left[  U\right]  }\left\langle f,h_{I;\kappa
}\right\rangle h_{I;\kappa}\ \text{and }g=\sum_{K\in\mathcal{G}}\left\langle
g,h_{K;\kappa}\right\rangle h_{K;\kappa}\ ,
\]
and so%
\begin{align}
\left\langle T_{\mathcal{G}}^{\eta}\mathbf{1}_{U}f,\mathbf{1}_{2^{n}%
U\setminus2^{n-1}U}g\right\rangle  & =\left\langle T_{\mathcal{G}}^{\eta}%
\sum_{I\in\mathcal{G}\left[  U\right]  }\left\langle f,h_{I;\kappa
}\right\rangle h_{I;\kappa},\mathbf{1}_{2^{n}U\setminus2^{n-1}U}\sum
_{K\in\mathcal{G}}\left\langle g,h_{K;\kappa}\right\rangle h_{K;\kappa
}\right\rangle \label{<Tf,g>}\\
& =\left\langle T_{\mathcal{G}}^{\eta}\sum_{I\in\mathcal{G}\left[  U\right]
}\left\langle f,h_{I;\kappa}\right\rangle h_{I;\kappa},\sum_{K\in
\mathcal{G}\left[  2^{n}U\setminus2^{n-1}U\right]  }\left\langle
g,h_{K;\kappa}\right\rangle h_{K;\kappa}\right\rangle \nonumber\\
& +\left\langle T_{\mathcal{G}}^{\eta}\sum_{I\in\mathcal{G}\left[  U\right]
}\left\langle f,h_{I;\kappa}\right\rangle h_{I;\kappa},\sum_{K\in
\mathcal{G}\setminus\mathcal{G}\left[  2^{n}U\setminus2^{n-1}U\right]
}\left\langle g,h_{K;\kappa}\right\rangle h_{K;\kappa}\right\rangle .\nonumber
\end{align}


For the inner product in the second line we have%
\begin{equation}
\left\langle T_{\mathcal{G}}^{\eta}\sum_{I\in\mathcal{G}\left[  U\right]
}\left\langle f,h_{I;\kappa}\right\rangle h_{I;\kappa},\sum_{K\in
\mathcal{G}\left[  2^{n}U\setminus2^{n-1}U\right]  }\left\langle
g,h_{K;\kappa}\right\rangle h_{K;\kappa}\right\rangle =\sum_{I\in
\mathcal{G}\left[  U\right]  }\sum_{J\in\mathcal{G}\left[  2^{n}%
U\setminus2^{n-1}U\right]  }\left\langle f,h_{I;\kappa}\right\rangle
\left\langle g,h_{J;\kappa}\right\rangle \left\langle T_{\mathcal{G}}^{\eta
}h_{I;\kappa},h_{J;\kappa}\right\rangle ,\label{well loc}%
\end{equation}
for all $I,J\in\mathcal{G}$, where we have replaced the dummy variable $K$ by
$J$, so as to avoid confusion when applying Lemma \ref{lem special},%
\begin{align*}
& \left\vert \left\langle T_{\mathcal{G}}^{\eta}\sum_{I\in\mathcal{G}\left[
U\right]  }\left\langle f,h_{I;\kappa}\right\rangle h_{I;\kappa},\sum
_{K\in\mathcal{G}\left[  2^{n}U\setminus2^{n-1}U\right]  }\left\langle
g,h_{K;\kappa}\right\rangle h_{K;\kappa}\right\rangle \right\vert \\
& \leq\sum_{I\in\mathcal{G}\left[  U\right]  }\sum_{J\in\mathcal{G}\left[
2^{n}U\setminus2^{n-1}U\right]  }\left\vert \left\langle f,h_{I;\kappa
}\right\rangle \right\vert \left\vert \left\langle g,h_{J;\kappa}\right\rangle
\right\vert \left\vert \left\langle T_{\mathcal{G}}^{\eta}h_{I;\kappa
},h_{J;\kappa}\right\rangle \right\vert \\
& \leq\sum_{I\in\mathcal{G}\left[  U\right]  }\sum_{J\in\mathcal{G}\left[
2^{n}U\setminus2^{n-1}U\right]  }\left\vert \left\langle f,h_{I;\kappa
}\right\rangle \right\vert \left\vert \left\langle g,h_{J;\kappa}\right\rangle
\right\vert C_{d-1}\frac{1}{\eta^{\kappa}}\left(  \frac{\ell\left(  I\right)
}{2^{n}}\right)  ^{\kappa+\frac{d-1}{2}}\left(  \frac{\ell\left(  J\right)
}{2^{n}}\right)  ^{\kappa+\frac{d-1}{2}},
\end{align*}
where we have used that $T_{\mathcal{G}}^{\eta}$ is self-adjoint.

Now we see that the inner product in the middle line of (\ref{<Tf,g>}) is
controlled by%
\begin{align*}
& \left\vert \left\langle T_{\mathcal{G}}^{\eta}\sum_{I\in\mathcal{G}\left[
U\right]  }\left\langle f,h_{I;\kappa}\right\rangle h_{I;\kappa},\sum
_{J\in\mathcal{G}\left[  2^{n}U\setminus2^{n-1}U\right]  }\left\langle
g,h_{J;\kappa}\right\rangle h_{J;\kappa}\right\rangle \right\vert \\
& \leq\sum_{I\in\mathcal{G}\left[  U\right]  }\sum_{J\in\mathcal{G}\left[
2^{n}U\setminus2^{n-1}U\right]  }\left\vert \left\langle f,h_{I;\kappa
}\right\rangle \left\langle g,h_{J;\kappa}\right\rangle \right\vert
C_{d-1}\frac{1}{\eta^{\kappa}}\left(  \frac{\ell\left(  I\right)  }{2^{n}%
}\right)  ^{\kappa+\frac{d-1}{2}}\left(  \frac{\ell\left(  J\right)  }{2^{n}%
}\right)  ^{\kappa+\frac{d-1}{2}},
\end{align*}
which is equal to%
\begin{align*}
& \sum_{I\in\mathcal{G}\left[  U\right]  }\sum_{J\in\mathcal{G}\left[
2^{n}U\setminus2^{n-1}U\right]  }\frac{\left(  \int\left\vert \bigtriangleup
_{I;\kappa}f\right\vert \right)  \left(  \int\left\vert \bigtriangleup
_{J;\kappa}g\right\vert \right)  }{\ell\left(  I\right)  ^{-\frac{d-1}{2}}%
\ell\left(  J\right)  ^{-\frac{d-1}{2}}}C_{d-1}\frac{1}{\eta^{\kappa}}\left(
\frac{\ell\left(  I\right)  }{2^{n}}\right)  ^{\kappa+\frac{d-1}{2}}\left(
\frac{\ell\left(  J\right)  }{2^{n}}\right)  ^{\kappa+\frac{d-1}{2}}\\
& =\left(  \int\sum_{I\in\mathcal{G}\left[  U\right]  }\frac{\left\vert
\bigtriangleup_{I;\kappa}f\right\vert }{\ell\left(  I\right)  ^{-\frac{d-1}%
{2}}}\left(  \frac{\ell\left(  I\right)  }{2^{n}}\right)  ^{\kappa+\frac
{d-1}{2}}\right)  \left(  \int\sum_{J\in\mathcal{G}\left[  2^{n}%
U\setminus2^{n-1}U\right]  }\frac{\left\vert \bigtriangleup_{J;\kappa
}g\right\vert }{\ell\left(  J\right)  ^{-\frac{d-1}{2}}}C_{d-1}\frac{1}%
{\eta^{\kappa}}\left(  \frac{\ell\left(  J\right)  }{2^{n}}\right)
^{\kappa+\frac{d-1}{2}}\right)  \equiv AB.
\end{align*}


The first factor $A$ satisfies,%
\begin{align*}
A  & =\int\sum_{I\in\mathcal{G}\left[  U\right]  }\frac{\left\vert
\bigtriangleup_{I;\kappa}f\right\vert }{\ell\left(  I\right)  ^{-\frac{d-1}%
{2}}}\left(  \frac{\ell\left(  I\right)  }{2^{n}}\right)  ^{\kappa+\frac
{d-1}{2}}\\
& \leq\int_{\left(  1+\eta\right)  U}\left(  \sum_{I\in\mathcal{G}\left[
U\right]  }\left\vert \bigtriangleup_{I;\kappa}f\right\vert ^{2}\right)
^{\frac{1}{2}}\left(  \sum_{I\in\mathcal{G}\left[  U\right]  }\left(  \frac
{1}{\ell\left(  I\right)  ^{-\frac{d-1}{2}}}\left(  \frac{\ell\left(
I\right)  }{2^{n}}\right)  ^{\kappa+\frac{d-1}{2}}\right)  ^{2}\right)
^{\frac{1}{2}}\\
& \leq\left(  \int\left(  \sum_{I\in\mathcal{G}\left[  U\right]  }\left\vert
\bigtriangleup_{I;\kappa}f\right\vert ^{2}\right)  ^{\frac{p}{2}}\right)
^{\frac{1}{p}}\left(  \sum_{I\in\mathcal{G}\left[  U\right]  }\left(  \frac
{1}{\ell\left(  I\right)  ^{-\frac{d-1}{2}}}\left(  \frac{\ell\left(
I\right)  }{2^{n}}\right)  ^{\kappa+\frac{d-1}{2}}\right)  ^{2}\right)
^{\frac{1}{2}}\\
& \leq\left\Vert f\right\Vert _{L^{p}}\left(  \sum_{I\in\mathcal{G}\left[
U\right]  }\left(  \frac{1}{\ell\left(  I\right)  ^{-\frac{d-1}{2}}}\left(
\frac{\ell\left(  I\right)  }{2^{n}}\right)  ^{\kappa+\frac{d-1}{2}}\right)
^{2}\right)  ^{\frac{1}{2}}\lesssim2^{-n\left(  2\kappa+d-1\right)
}\left\Vert f\right\Vert _{L^{p}}%
\end{align*}
by the square function theorem for Alpert wavelets, and since%
\begin{align*}
& \sum_{I\in\mathcal{G}\left[  U\right]  }\left(  \frac{1}{\ell\left(
I\right)  ^{-\frac{d-1}{2}}}\left(  \frac{\ell\left(  I\right)  }{2^{n}%
}\right)  ^{\kappa+\frac{d-1}{2}}\right)  ^{2}=\sum_{I\in\mathcal{G}\left[
U\right]  }\frac{1}{\ell\left(  I\right)  ^{-\left(  d-1\right)  }}\left(
\frac{\ell\left(  I\right)  }{2^{n}}\right)  ^{2\kappa+d-1}\\
& =\sum_{m=0}^{\infty}\frac{2^{m\left(  d-1\right)  -\left(  2m\kappa
+2m\left(  d-1\right)  \right)  }}{2^{n\left(  2\kappa+d-1\right)  }}%
=\sum_{m=0}^{\infty}\frac{2^{-m\left(  2\kappa+\left(  d-1\right)  \right)  }%
}{2^{n\left(  2\kappa+d-1\right)  }}=C2^{-n\left(  2\kappa+d-1\right)  }.
\end{align*}
The second factor is handled similarly to obtain%
\[
\left\vert \int\sum_{J\in\mathcal{G}\left[  2^{n}U\setminus2^{n-1}U\right]
}\frac{\left\vert \bigtriangleup_{J;\kappa}g\right\vert }{\ell\left(
J\right)  ^{-\frac{d-1}{2}}}C_{d-1}\frac{1}{\eta^{\kappa}}\left(  \frac
{\ell\left(  J\right)  }{2^{n}}\right)  ^{\kappa+\frac{d-1}{2}}\right\vert
\lesssim2^{-n\left(  2\kappa+d-1\right)  }\left\Vert f\right\Vert _{L^{p}}.
\]


It remains to control the term in the third line of (\ref{<Tf,g>}), namely%
\begin{align*}
& \left\langle T_{\mathcal{G}}^{\eta}\sum_{I\in\mathcal{G}\left[  U\right]
}\left\langle f,h_{I;\kappa}\right\rangle h_{I;\kappa},\sum_{K\in
\mathcal{G}\setminus\mathcal{G}\left[  2^{n}U\setminus2^{n-1}U\right]
}\left\langle g,h_{K;\kappa}\right\rangle h_{K;\kappa}\right\rangle \\
& =\sum_{I\in\mathcal{G}\left[  U\right]  }\sum_{K\in\mathcal{G}%
\setminus\mathcal{G}\left[  2^{n}U\setminus2^{n-1}U\right]  }\left\langle
f,h_{I;\kappa}\right\rangle \left\langle g,h_{K;\kappa}\right\rangle
\left\langle T_{\mathcal{G}}^{\eta}h_{I;\kappa},h_{K;\kappa}\right\rangle \\
& =\sum_{I\in\mathcal{G}\left[  U\right]  }\sum_{J\in\mathcal{G}%
\setminus\mathcal{G}\left[  2^{n}U\setminus2^{n-1}U\right]  }\left\langle
f,h_{I;\kappa}\right\rangle \left\langle g,h_{J;\kappa}\right\rangle
\left\langle T_{\mathcal{G}}^{\eta}h_{I;\kappa},h_{J;\kappa}\right\rangle ,
\end{align*}
where $\operatorname*{Supp}g\subset2^{n}U\setminus2^{n-1}U$, and we gave again
replaced the dummy variable $K$ with $J$. Then $\left\langle g,h_{J;\kappa
}\right\rangle \neq0$ implies $J\cap2^{n}U\setminus2^{n-1}U\neq\emptyset$, and
we see that such $J\,$must be far from $I$. But when $J$ is far from $I$, this
sum can be handled by the same techniques we just used for the second line in
(\ref{<Tf,g>}). This completes the proof of (\ref{lin d'}) and Corollary
\ref{Cor LSST}.
\end{proof}

\end{document}
